\chapter{Scope of this edition}\label{ch:limits}

This reference covers message framing, the bulk-dump transfer protocol, the parameter
address space, the universal messages, error reporting and the differences between the two
models. The following are not covered.

\begin{description}
\item[What the parameters do.] Chapter~\ref{ch:params} enumerates all 108 addresses of
  the \bytes{00} area with their name, length, bit mask, accepted values and position in a
  dump. It does not explain what each setting does to the sound; the user manuals do that.
\item[The effect parameters.] The twenty \textsc{value} bytes under each of
  \textsc{effect1}, \textsc{effect2} and \textsc{reverb} mean different things for each
  effect type. The published guide gives them, effect by effect, and this edition does not
  reproduce them.
\item[Block contents.] Chapter~\ref{ch:blocks} gives each block's signature and where its
  names are, but not how a stored sound, combination or sequence is arranged inside one.
  Chapter~\ref{ch:sound} lays out the sound being \emph{edited}, which is a different
  thing from the sound memory a bulk dump carries.
\item[The SEQUENCER category, against hardware.] Its protocol is described, but the
  capture used to check everything else came from a rack, which cannot send it. See
  chapter~\ref{ch:overview}.
\item[Operating system version 1.] Everything here describes version 2.0.
\end{description}
